% Chapter 4 — includable by main.tex
\chapter{Ekvivalentnost klasa regularnih jezika}
\label{chap:4}

\begin{quote}
\itshape
U dosad uveli nekoliko načina opisivanja regularnih jezika, a ovo
poglavlje zatvara krug — pokazat ćemo ekvivalentnost različitih formalizama.
\end{quote}

\newpage

% ================================================================
\section{Krug koji treba zatvoriti}
% ================================================================

Dosad smo uveli nekoliko načina opisivanja regularnih jezika:
\begin{itemize}
  \item \textbf{RI} --- regularni izrazi (poglavlje 2)
  \item \textbf{KA} --- deterministički konačni automati (poglavlje 3)
  \item \textbf{NKA} --- nedeterministički konačni automati (poglavlje 3)
  \item \textbf{DLG} --- desnolinearne gramatike (ovo poglavlje)
  \item \textbf{LLG} --- lijevolinearne gramatike (ovo poglavlje)
\end{itemize}

Svaki od ovih formalizma definira neku klasu jezika. Ključni teorem
ovog poglavlja kaže da su sve te klase iste:
\[
  \mathrm{Reg_{KA}} = \mathrm{Reg_{NKA}} = \mathrm{Reg_{RI}}
  = \mathrm{Reg_{DLG}} = \mathrm{Reg_{LLG}}.
\]
Tu klasu jednostavno zovemo $\mathrm{Reg}$ --- \textbf{klasa regularnih
jezika}.

Dokaz ide ``u krug'': pokazujemo $\mathrm{RI} \to \mathrm{NKA} \to
\mathrm{KA} \to \mathrm{DLG} \to \mathrm{RI}$, a $\mathrm{LLG}$
uključujemo pomoću operacije reverza.

% ================================================================
\section{Desnolinearne gramatike (DLG)}
% ================================================================

\begin{definicija}[Gramatika]
\textbf{Gramatika} nad abecedom $\Sig$ je transformacijski sustav
$G = (V, \Sig, \to, S)$, gdje je:
\begin{itemize}
  \item $V$ --- konačni skup \textbf{varijabli} (disjunktan od $\Sig$),
  \item $Y := \Sig \cup V$ --- skup svih \textbf{simbola},
  \item $S \in V$ --- \textbf{početna varijabla},
  \item ${\to} \subseteq Y^+ \times Y^*$ --- konačna relacija
        \textbf{pravila} (na lijevoj strani barem jedna varijabla).
\end{itemize}
Pisemo $u \Rightarrow v$ ako postoji pravilo $s \to t$ i rijeci $a, b$
takve da je $u = asb$ i $v = atb$. Gramatika \textbf{generira} jezik
\[
  L(G) := \{w \in \Sig^* : S \Rightarrow^* w\}.
\]
\end{definicija}

Ovo je opća definicija gramatike. Za regularne jezike dovoljne su nam
gramatike s \emph{jako} ograničenim pravilima:

\begin{definicija}[Desnolinearna gramatika]
Gramatika $G$ je \textbf{desnolinearna} (DLG) ako su sva njena pravila
jednog od dva oblika:
\[
  A \to \alpha B \qquad \text{ili} \qquad A \to \eps,
\]
gdje su $A, B \in V$ varijable, a $\alpha \in \Sig$ znak.
\end{definicija}

Intuicija: svako pravilo ili ``pojede'' jedan znak i zamijeni varijablu
s novom varijablom (koja će pojesti sljedeći), ili varijabla nestane
(kraj izvoda). Varijabla je uvijek na \emph{desnom} kraju međurječita
--- otuda ime ``desnolinearna''.

\begin{primjer}
DLG s varijablama $\{S, A\}$ i pravilima
$S \to aA$, $A \to aA$, $A \to bA$, $A \to \eps$
generira jezik svih rijeci koje pocinju s $a$ i zatim mogu imati bilo
sto. Regularni izraz: $a(a \cup b)^*$.
\end{primjer}

% ================================================================
\section{Pretvorba KA $\to$ DLG}
% ================================================================

Svaki KA mozemo pretvoriti u ekvivalentnu DLG na prirodan nacin:
stanja postaju varijable, prijelazi postaju pravila.

\begin{propozicija}
Svaki KA-regularni jezik je DLG-regularan.
\end{propozicija}

\textit{Dokaz (konstrukcija).}
Neka je $M = (Q, \Sig, \delta, q_0, F)$ KA koji prepoznaje jezik $L$.
Konstruiramo DLG $G$ ovako:
\begin{itemize}
  \item Za svako stanje $q \in Q$ uvedemo varijablu $V_q$.
  \item Početna varijabla: $S := V_{q_0}$.
  \item Za svaki prijelaz $\delta(q, \alpha) = p$ dodamo pravilo
        $V_q \to \alpha V_p$.
  \item Za svako završno stanje $q \in F$ dodamo pravilo $V_q \to \eps$.
\end{itemize}
Moze se pokazati indukcijom da za svako stanje $q \in Q$ vrijedi
\[
  L(Q, \Sig, \delta, q, F) = L(V, \Sig, \to, V_q),
\]
tj.\ varijabla $V_q$ generira tocno rijeci koje KA prihvaca krenuvsi
iz stanja $q$. Za $q := q_0$ dobijemo $L(M) = L(G)$. \qed

\begin{primjer}
KA za rijeci koje zavrsavaju na $ab$ iz poglavlja 3 pretvara se u DLG:
\begin{align*}
  V_{q_0} &\to a V_{q_a} \mid b V_{q_0} \\
  V_{q_a} &\to a V_{q_a} \mid b V_{q_{ab}} \\
  V_{q_{ab}} &\to a V_{q_a} \mid b V_{q_0} \mid \eps
\end{align*}
Ovo generira točno riječi koje završavaju na $ab$.
\end{primjer}

% ================================================================
\section{Gramaticki sustav i Ardenovo pravilo}
% ================================================================

Za smjer $\mathrm{DLG} \to \mathrm{RI}$ koristimo tehniku
\textbf{gramatickih jednadžbi}.

\subsection{Gramaticki sustav}

Svaka DLG zadaje \emph{sustav jednadžbi} nad jezicima: za svaku
varijablu $A$ skupimo sva njena pravila i zapisemo jednadžbu:
\[
  A = \alpha_1 B_1 \cup \alpha_2 B_2 \cup \cdots \cup \alpha_k B_k
  \;\cup\; (\eps \text{ ako ima pravilo } A \to \eps).
\]
Ako $A$ nema nijednog pravila, dodamo $A = \emptyset$.

Supstitucijom (uvrstavanjem jednadžbi jedne u drugu) pokusavamo
eliminirati varijable i doci do regularnog izraza za $S$.

\subsection{Problem: rekurzija}

Supstitucija funkcionira kada nema rekurzije. Ali DLG cesto ima
rekurzivna pravila poput $A = aA \cup b$, gdje se $A$ pojavljuje
s obje strane. Sto tada?

\begin{propozicija}[Ardenovo pravilo]
Jednadžba $X = AX \cup B$ ima rjesenje $X = A^* B$.
Ako je $A$ pozitivan, to je \emph{jedino} rjesenje.
Inace, to je \emph{najmanji} skup koji zadovoljava jednadžbu.
\end{propozicija}

\textit{Provjera.} Uvrstimo $X = A^*B$ u jednadžbu:
\[
  A(A^*B) \cup B = (AA^*)B \cup B = A^+B \cup \eps B = (A^+ \cup \eps)B = A^*B. \; \checkmark
\]

Ardenovo pravilo je analogon formule za geometrijski red:
$a \cdot r + b = r$ ima rjesenje $r = b/(1-a)$, tj.\ $r = a^* b$
u jezicnoj algebri.

\begin{primjer}[Pretvorba DLG $\to$ RI]
Pretvorimo DLG s pravilima:
\[
  S \to aA \mid bB, \quad
  A \to aA \mid bA \mid aN, \quad
  B \to aB \mid bB \mid bN, \quad
  N \to \eps.
\]
Gramaticki sustav ($N = \eps$ odmah uvrstimo):
\begin{align*}
  S &= aA \cup bB \\
  A &= aA \cup bA \cup a = (a \cup b)A \cup a
\end{align*}
Po Ardenovom pravilu ($a \cup b$ je pozitivan):
$A = (a \cup b)^* a$.

Analogno: $B = (a \cup b)^* b$.

Uvrstimo u $S$:
\[
  S = a(a \cup b)^*a \cup b(a \cup b)^*b.
\]
To su sve riječi koje počinju i završavaju istim slovom! \checkmark
\end{primjer}

% ================================================================
\section{RI $\to$ NKA: Thompsonova konstrukcija}
% ================================================================

Sad zatvaramo krug s druge strane: svaki regularni izraz mozemo
pretvoriti u NKA. Dokaz ide indukcijom po izgradnji izraza.

\begin{propozicija}
Svaki RI-regularni jezik je NKA-regularan.
\end{propozicija}

\textit{Dokaz (Thompsonova konstrukcija).}
Za svaki inicijalni jezik konstruiramo trivijalni NKA:
\begin{itemize}
  \item $\emptyset$: jedno stanje, nema prijelaza, nema zavrsnih stanja.
  \item $\eps$: jedno stanje, ono je i početno i završno, nema prijelaza.
  \item $\{\alpha\}$: dva stanja, prijelaz $\alpha$ od početnog do završnog.
\end{itemize}
Za slozenije izraze koristimo konstrukcije iz poglavlja 3
(unija s $\eps$-prijelazima, konkatenacija, zvijezda). \qed

% ================================================================
\section{Lijevolinearne gramatike i reverzi}
% ================================================================

\begin{definicija}[Lijevolinearna gramatika]
Gramatika je \textbf{lijevolinearna} (LLG) ako su sva njena pravila
oblika:
\[
  A \to B\alpha \qquad \text{ili} \qquad A \to \eps,
\]
dakle varijabla je na \emph{lijevom} kraju desne strane pravila.
\end{definicija}

LLG generiraju iste jezike kao i DLG (a time i sve ostale klase),
ali to nije odmah ocito. Dokazujemo to pomocu \textbf{reverza}.

\begin{definicija}[Reverz]
\textbf{Reverz} rijeci $w = \alpha_1 \alpha_2 \cdots \alpha_n$ je
$w^R = \alpha_n \cdots \alpha_2 \alpha_1$.
\textbf{Reverz jezika} je $L^R = \{w^R : w \in L\}$.
\end{definicija}

Za regularne izraze vrijede jednostavne formule:
\[
  \emptyset^R = \emptyset, \quad \eps^R = \eps, \quad \alpha^R = \alpha,
\]
\[
  (LM)^R = M^R L^R, \quad (L \cup M)^R = L^R \cup M^R, \quad (L^*)^R = (L^R)^*.
\]

Iz ovih formula indukcijom slijedi:

\begin{propozicija}
Reverz svakog RI-regularnog jezika je ponovo RI-regularan.
Dakle, klasa $\mathrm{Reg_{RI}}$ zatvorena je na reverz.
\end{propozicija}

Sada, DLG koja generira jezik $L$ lako se pretvori u LLG koja generira
$L^R$ (samo okrenimo redoslijed u pravilima). Kako je $\mathrm{Reg_{RI}}$
zatvorena na reverz, zakljucujemo:
\[
  \mathrm{Reg_{LLG}} = (\mathrm{Reg_{DLG}})^R = (\mathrm{Reg_{RI}})^R = \mathrm{Reg_{RI}}.
\]

% ================================================================
\section{Zakljucak: sve je isto}
% ================================================================

Evo kompletnog kruga pretvorbi:

\begin{center}
\renewcommand{\arraystretch}{1.6}
\begin{tabular}{lll}
\hline
\textbf{Od} & \textbf{Do} & \textbf{Metoda} \\
\hline
RI & NKA & Thompsonova konstrukcija (indukcija po izrazu) \\
NKA & KA & Partitivna konstrukcija \\
KA & DLG & Stanja $\leftrightarrow$ varijable, prijelazi $\leftrightarrow$ pravila \\
DLG & RI & Gramaticki sustav + Ardenovo pravilo \\
DLG & LLG & Reverz (okretanje pravila) \\
\hline
\end{tabular}
\end{center}

\begin{kljucno}
Sve ove pretvorbe su algoritamske --- dajte mi jedno, dat cu vam
bilo koje drugo. Svi formalizmi su jednako mocni; razlika je samo
u pogodnosti za razlicite svrhe. RI je dobar za opis, KA za
implementaciju, DLG za vezu s gramatikama.
\end{kljucno}

\begin{zadatak}
\begin{enumerate}[label=(\alph*), itemsep=4pt]
  \item Pretvorite regularni izraz $a^*(b \cup c)^+$ u NKA
        Thompsonovom konstrukcijom.
  \item Primijenite partitivnu konstrukciju na rezultat.
  \item Pretvorite KA koji ste dobili u DLG.
  \item Rijesite gramaticki sustav te DLG i dobijte regularni izraz.
        Usporedite s polaznim izrazom.
  \item Napisite LLG ekvivalentnu DLG iz (c).
\end{enumerate}
\end{zadatak}

% ================================================================
\section*{Sazetak}
% ================================================================

Pokazali smo da pet naizgled razlicitih formalizma opisuju istu klasu:

\medskip
\noindent
\begin{tabular}{ll}
  $\mathrm{Reg_{KA}}$ & jezici KA \\
  $\mathrm{Reg_{NKA}}$ & jezici NKA \\
  $\mathrm{Reg_{RI}}$ & regularni izrazi \\
  $\mathrm{Reg_{DLG}}$ & desnolinearne gramatike \\
  $\mathrm{Reg_{LLG}}$ & lijevolinearne gramatike \\
\end{tabular}

\medskip
Sve su iste klase. Tu klasu zovemo $\mathrm{Reg}$ --- \textbf{regularni
jezici}. U sljedecem poglavlju prelazimo na bezkontekstne jezike ---
prvu pravu nadklasu.

